\def\stat{horoshilov}





\def\tit{ТОКЕНИЗАЦИЯ ТЕКСТОВ\\ НА~ОСНОВЕ МЕТОДА ФУНКЦИОНАЛЬНЫХ ШАБЛОНОВ}





\def\titkol{Токенизация текстов на~основе метода функциональных шаблонов}





\def\aut{Ю.\,В.~Никитин$^1$, А.\,А.~Хорошилов$^2$, А.\,Е.~Макарова$^3$}





\def\autkol{Ю.\,В.~Никитин, А.\,А.~Хорошилов, А.\,Е.~Макарова}





\titel{\tit}{\aut}{\autkol}{\titkol}





\index{Никитин Ю.\,В.}


\index{Хорошилов А.\,А.}


\index{Макарова А.\,Е.}


\index{Nikitin Yu.\,V.}


\index{Khoroshilov A.\,A.}


\index{Makarova A.\,E.}





%{\renewcommand{\thefootnote}{\fnsymbol{footnote}}


%\footnotetext[1]


%{Работа выполнена в~МСЦ РАН в~рамках государственного задания по теме FNEF-2022-0014.}} 





\renewcommand{\thefootnote}{\arabic{footnote}}


\footnotetext[1]{Федеральный исследовательский центр <<Информатика и~управление>> Российской академии наук, 


\mbox{yuri.v.nikitin@gmail.com}}


\footnotetext[2]{Московский авиационный институт; Федеральный исследовательский центр <<Информатика и~управление>> 


Российской академии наук; 27~Центральный научно-исследовательский институт Министерства обороны 


Российской Федерации, \mbox{khoroshilov@mail.ru}}


\footnotetext[3]{АО <<НПК ``ВТ и~СС''>>, anna20497@list.ru}





 \label{st\stat}





 \thispagestyle{headings}


 


 \vspace*{-12pt} 


 


  


  


   


  \Abst{Предлагается новый метод токенизации текста, базирующийся на применении 


обобщенных функциональных шаблонов. В~основу метода положена классификация 


символов Юникода, учитывающая их роль в~формировании элементарных элементов текста 


(сегментов), и~классификация сформированных сегментов по типу их графематических 


классов. Особенность метода заключается в~использовании при формировании шаблона 


интервалов последовательности символов Юникода. Основное достоинство метода~--- 


возможность токенизации сложных информационных объектов (номера, географические 


координаты, наименования артикулов инженерных изделий и~т.\,п.), возможность получения 


детальной классификации токенов на стадии их формирования, возможность включения 


и~выключения токенизации определенного вида токенов, а также добавления новых шаблонов 


по образцу текста для дообучения системы.}


  


  \KW{токенизация; сегментация текста; функциональные шаб\-ло\-ны; графематический 


анализ; компьютерная лингвистика; автоматическая обработка текстов}





\DOI{10.14357\!/\!08696527220406}  





\vskip 10pt plus 9pt minus 6pt





 \thispagestyle{headings}


 


 \vspace*{-18pt}


  


\section{Введение}





\vspace*{-2pt}


  


  В процессе автоматической обработки текстов (АОТ) на начальном этапе 


необходимо преобразовать символьное представление текста 


в~последовательность его текстовых структурных объектов. Такое 


преобразование возможно осуществить только на основе закономерностей 


установления границ между элементами текста. В~рамках процесса АОТ это 


преобразование называется токенизацией. 


  


  В настоящее время в~системах обработки текстовой информации 


используются различные методы токенизации: методы, базирующиеся на 


регулярных выражениях~[1]; методы, основанные на правилах~[2]; 


нейросетевые методы~[3]. Методы, основанные на регулярных выражениях 


и~на правилах, не могут учесть все разнообразие сложных случаев, а~нейросети 


реализуют только те ситуации, которые встречались в~обучающей выборке.


  


  Структурные элементы текста, выделяемые средствами токенизации, 


относятся к~различным информационным объектам. К~текстовым объектам 


относятся слова (простые, сложные, аббревиатуры, хештеги, слова, содержащие 


в~своем составе числа, и~т.\,п.), знаки пунктуации (в~том числе кавычки 


и~скобки), циф\-ро\-вые последовательности (положительные, отрицательные 


и~дробные числа, проценты, даты, числа в~экспоненциальной записи или 


с~разделением разрядов и~т.\,п.), знаки обозначения (амперсанд, параграф, 


номер, копирайт, обозначения валют, символы математических операций 


и~т.\,п.). 





\section{Основная идея предлагаемого метода}


  


  Процесс токенизации текстов состоит из следующей последовательности 


процедур:


  \begin{itemize}


\item сегментация~--- деление текста на сегменты по графематическим 


классам символов;


\item формально-логический контроль~--- замена омоглифов и~букв~<<ё>>, исправление частых ошибок в~словах;


\item токенизация в~узком смысле~--- объединение сегментов в~токены;


\item детальная классификация токенов по типу их графематических 


классов.


\end{itemize}





  \textbf{Сегментация.} Под сегментацией понимается деление текста на 


элементарные сегменты, состоящие из символов с~однородными 


графематическими классами. Под однородностью графематических классов 


понимается возможность их сочетания в~одном сегменте (настраиваемый 


параметр). Примерами сегментации, при которой графические отображения 


символов с~однородными классами собираются в~единый сегмент, могут 


служить следующие случаи: L и~l (основная латиница); K и~k (кириллица); G и~g 


(греческий).


  


  Последовательности разделителей элементов текста (пробелы, табуляции, 


переводы строки, разделители звездочками, отчеркивания и~т.\,п.)\ в~процессе 


сегментации рассматриваются отдельно. Такие последовательности между 


токенами называем делимитерами. Они делятся на два вида:


  \begin{enumerate}[(1)]


\item пробельные~--- состоят из любого числа вышеперечисленных 


разделителей текста, кроме перевода строки;


\item концы строк~--- состоят из любого числа разделителей, включая хотя 


бы один перевод строки.


\end{enumerate}





  Все пробельные делимитеры, независимо от их числа, заменяются 


одиночными пробелами. Аналогично все делимитеры концов строк заменяются 


на один перевод строки.


  


  Таким образом, с~помощью преобразования делимитеров 


к~унифицированному разделителю (про\-бел\!/\!пе\-ре\-вод строки) текст 


очищается для упрощения дальнейшего анализа.


  


  Результатом сегментации является кортеж (упорядоченный массив) 


элементарных сегментов текста, состоящих из однородных символов текста 


или делимитеров.


  


  \textbf{Формально-логический контроль.} На данном этапе в~словах 


проводится замена буквы <<ё>> на <<е>> и~исправление омоглифов (букв, 


имеющих совпадающее написание в~разных алфавитах) на буквы <<родного>> 


алфавита. После этого для слов с~обнаруженными отклонениями проводятся 


замены символов и~измененный текст сегмента записывается в~эталонное 


значение контента. При этом исходный контент сегмента сохраняется. При 


дальнейшем разборе для слов, имеющих эталонный контент, в~качестве их 


текстового значения рассматривается именно эталон.


  


  \textbf{Токенизация.} Под токенизацией (в~узком смысле, в~рамках карты 


процесса АОТ) понимается процесс объединения элементарных сегментов 


текста в~составной токен.


  


  Ввиду того что алфавитные символы, цифры и~знаки (в~том числе дефисы 


и~апострофы, играющие в~словах роль буквы) имеют неоднородные 


графематические классы, они попадают в~разные сегменты, поэтому для 


сложных последовательностей (см.\ таблицу) требуется дополнительная 


<<сборка>> в~один элемент (токен).


  \begin{table}[b]\small


\begin{center}


\begin{tabular}{|l|l|}


\multicolumn{2}{c}{Примеры сложных токенов}\\


\multicolumn{2}{c}{| }\\[-6pt]


\hline


\multicolumn{1}{|c|}{\textbf{Обозначение}}&\multicolumn{1}{c|}{\textbf{Название}}\\


\hline


8-800-000-00-00&Телефон\\


\hline


&\\[-9pt]


55$^\circ$45$^\prime$21$^{\prime\prime}$ с.~ш. 37$^\circ$37$^\prime$04$^{\prime\prime}$ 


в.~д.&Географические координаты\\


\hline


1.022E23&Экспоненциальная запись числа\\


\hline


$+34$~$^\circ$C&Температура\\


\hline


postbox@mail.ru&Электронная почта\\


\hline


{\sf www.yandex.ru}&Адрес сайта\\


\hline


\#москва&Хештег\\


\hline


@kremlinmuseums&Аккаунт в~соцсети\\


\hline


01.01.2022&Дата\\


\hline


09:02&Время\\


\hline


&\\[-9pt]


:)&Смайл\\[1pt]


\hline


$10 - 7 = 3$&Математическое выражение\\


\hline


Вице-президент&Слова, написанные через дефис\\


\hline


Emias.info, Грамота.ру, Яндекс.Карты&Слова, содержащие в~своем составе точку\\


\hline


Москва~--- столица &Слова, написанные через тире \\


\hline


Москва-92, Су-57&Слова, имеющие в~составе числительное\\


\hline


\tabcolsep=0pt\begin{tabular}{l}300-летие, 100-руб\-ле\-вый,\\   15-про\-цент\-ный\end{tabular}&


\tabcolsep=0pt\begin{tabular}{l}Существительные и~прилагательные\\ в~словесно-цифровой форме\end{tabular} \\


\hline


Д'Артаньян&Слова, содержащие апостроф \\


\hline


\end{tabular}


\end{center}


\end{table}


  Такие конструкции представляют собой составные объединения 


элементарных сегментов. 





Каждый сегмент можно описать простейшей 


формулой на основе его графематического класса и~длины (по следующему 


принципу: <<одна буква>>, <<две цифры>>, <<один дефис>>, <<много 


букв>>, <<одна точка>>, <<много цифр>> и~т.\,п.). Тогда всю конструкцию 


можно описать составной формулой всех ее сегментов. Например: формула 


<<много букв, один дефис, много букв>> (9L~1-~9L) соответствует словам вида 


<<вице-президент>>; формула <<две цифры, одна точка, две цифры, одна 


точка, четыре цифры>> (2D~1.~2D~1.~4D) соответствует цифровому формату 


даты <<01.01.2022>>.


  Словарь формул составных конструкций формируется по корпусу текстов, 


а~объединение сегментов в~токены на этапе токенизации проводится уже по 


такому словарю.


  


  Аналогично в~единую составную конструкцию собираются последовательно 


стоящие делимитеры с~той лишь разницей, что составной сегмент не 


становится значимым токеном текста, а также унифицируется в~разделитель 


текста.


  


  Типы токенов (слово, чис\-ло\!/\!но\-мер, знак пунктуации, знак обозначения) и~тип 


<<делимитер>> определяются по словарю формул составных сегментов GA\_3, 


т.\,е.\ они заранее присваиваются формулам на этапе составления словаря из 


шаблонов текстовых фрагментов соответствующего вида.


  


  \textbf{Классификация токенов.} В~процессе токенизации проводится 


детальная классификация токенов: определяется типологический класс 


(простое слово, сложное слово, число, номер, знак препинания, знак 


обозначения), вид (например, дроб\-ное число, IP-ад\-рес, телефонный номер) 


и~графематический класс (классификация в~зависимости от  


бук\-вен\-но-зна\-ко\-во\-го состава и~регистра символов). Это исключает 


необходимость проведения ряда технологических операций на последующих 


стадиях обработки текстов. 





\section{Модель процесса}


  


  Сопоставление символов текста с~соответствующими графематическими 


классами описывается при помощи дискретной функции


  $$


  \mathrm{CharClass}\left( x_i\right)=y_i\,,


  $$


значения которой представляют собой элементы множества $\mathrm{Cls}$,


$$


\mathrm{Cls}= \{ y_i\}=\{\mathrm{L,l,K,k,D},\ldots\}.


$$


  


  Определенные заранее графематические классы занесены в~таблицу GA\_1:


  $$


  \mathrm{GA_1}= \begin{pmatrix}


  x_1 & y_1\\


  x_2 & y_2\\


  \cdots &\cdots\\


  x_s& y_s


  \end{pmatrix}.


  $$


Здесь $s=1,\ldots , e$, где $e$~--- число учитываемых символов исходного текста 


и~соответствующих им значений функции $\mathrm{CharClass}\,(x_i)$. 


  


  Сегменты segm$_*$ и~segm$_{**}$, формулы которых совпадают, считаются 


однородными:


  $$


  \mathrm{SymbForm}\left( \mathrm{segm}_* \right) =\mathrm{SymbForm}\left( 


\mathrm{segm}_{**}\right).


  $$


  


  Объединение элементарных сегментов текста в~составные (токены) будет 


выглядеть следующим образом:


  $$


  \mathrm{Token}= \left\{ \mathrm{segm}_1, \mathrm{segm}_2, \ldots , 


\mathrm{segm}_p\right\},


  $$


где $p$~--- число сегментов в~токене.


  


  Длины сегментов каждого графематического класса ранжируются по 


интервалам:


  $$


  \mathrm{Intervals}_{\mathrm{segm}}:


  \left[


  \begin{array}{rl}


  1\leq \mathrm{len}\,(\mathrm{segm})&\leq b_1\,;\\


  b_1+1\leq \mathrm{len}\,(\mathrm{segm})&\leq b_2\,;\\


  b_2+1\leq \mathrm{len}\,(\mathrm{segm})&\leq b_3\,;\\


  \cdots\\


  b_{n-1}+1\leq \mathrm{len}\,(\mathrm{segm})&\leq b_n\,,


  \end{array}


  \right.


  $$


где $n\hm= 1,\ldots , m$~--- число интервалов класса символа; $\mathrm{B}\hm= \{b_n\}$~--- границы интервалов. Полученные интервалы хранятся 


в~таблице GA\_2: 


$$


\mathrm{GA_2}= \begin{pmatrix}


y_1 & a_1^1 & b_1^1 & F_1^1 & F^{\prime1}_1\\


y_1 & \cdots&\cdots&\cdots&\cdots\\


y_1 & a_q^1 & b_q^1 & F_q^1 & F^{\prime1}_q\\


\cdots&\cdots&\cdots&\cdots&\cdots\\


y_j & a_z^j & b_z^j & F_z^j & F^{\prime j}_{z^\prime}\\


\cdots&\cdots&\cdots&\cdots&\cdots\\


y_s & a_n^s & b_n^s & F_n^s & F^{\prime s}_{n^\prime}


\end{pmatrix},


$$


где $a_n^s$~--- левая граница интервала для графематического класса 


символа~$y_s$;


$b_n^s$~--- правая граница интервала для графематического класса 


символа~$y_s$; 


$F_n^s$~--- формула графематического класса символа~$y_s$ в~интервале 


$[a_n^s; b_{n}^s]$; 


$F^{\prime s}_{n^\prime}$~--- символ формулы графематического класса символа~$y_s$ 


в~интервале $[a_n^s; b_n^s]$.


  


  Формула элементарного сегмента будет определена следующим образом:


  $$


  \mathrm{Form}\,(\mathrm{segm})= 


  \begin{cases}


  F_1\,, & \mathrm{segm}: 1\leq \mathrm{len}\,(\mathrm{segm})\leq b_1\,;\\


  F_2\,, & \mathrm{segm}: b_1+1\leq \mathrm{len}\,(\mathrm{segm})\leq b_2\,;\\


  F_3\,, & \mathrm{segm}: b_2+1\leq \mathrm{len}\,(\mathrm{segm})\leq b_3\,;\\


 & \hspace*{15mm}\cdots\\


  F_g\,, & \mathrm{segm}: b_{n-1}+1\leq \mathrm{len}\,(\mathrm{segm})\leq 


b_n\,.\\


  \end{cases}


  $$


Здесь $g=1,\ldots , c$, где $c$~--- число формул интервалов класса символа.   


  


  Тогда символ формулы элементарного сегмента будет записан как


  $$


  \mathrm{SymbForm}\,(\mathrm{segm})=F^\prime_{g^\prime} \vert 


\mathrm{Form}\,(\mathrm{segm})=F_g\,.


  $$


Здесь $g^\prime=1, \ldots , c^\prime$, где $c^\prime$~--- число символов формул 


интервалов.





  Объединение элементарных сегментов в~токен происходит по словарю 


шаблонов, составленному вручную с~учетом специфики конкретного текста. 


  


  В итоге составной сегмент будет представлять собой конкатенацию 


символьных строк формул сегментов $\mathrm{segm}_1, 


\mathrm{segm}_2,\ldots , \mathrm{segm}_k$:


  $$


  \mathrm{Comp\_segm}= \bigcup\limits_1^{\max} 


\mathrm{Form}\left(\mathrm{segm}_k\right)\,,


  $$


 а исходный текст будет определен как последовательность составных 


сегментов: 


  $$


  \mathrm{Text}=\left( \mathrm{segm}_1, \mathrm{segm}_2,\ldots , 


\mathrm{segm}_t \right),


  $$


где $t$~--- число составных сегментов в~тексте. 





\section{Алгоритм токенизации}


  





\noindent


\textbf{1.\ Замена символов исходного текста на символы 


графематических классов}


  


  Обозначение класса каждого символа определяется по таблице GA\_1.


  


  Исходный текст:


  


     \noindent


     {\sf  Вчера, 22.08.2020, экс-главу ЦРУ 8~часов допрашивали по делу 


о~вмешательстве в~выборы.}


  


  После замены:


  


     \noindent


    {\sf Kkkkl, DD.DD.DDDD, kkk-kkkkkKKKDkkkkkklkklkkkkkkkkklkkkkkkkkkkkkkkkk}


    


    \noindent


    {\sf kkkkkkk}.


      


      \smallskip





\noindent


\textbf{2.\ Поиск и~замена омоглифов и\!/\!или буквы <<ё>> при наличии 


таковых}


  


  Буква <<а>> в~слове <<Вчера>>, буквы <<о>> и~<<а>> в~слове 


<<допрашивали>>, а также <<е>> в~слове <<делу>> набраны на английской 


раскладке клавиатуры. После замены омоглифов соответствующие 


последовательности графематических классов изменятся:


  


  \noindent


      {\sf Kkkkk, DD.DD.DDDD, kkk-kkkkkKKKDkkkkkkkkkkkkkkkkkkkkkkkkkkkkkkkkkk}


      


      \noindent


      {\sf kkkkkkkkk}.


      


      


            \smallskip


      


\noindent


\textbf{3.\ Объединение в~один сегмент последовательностей 


с~однородными символами}


  


  Поскольку графематические классы символов <<{\sf K}>> и~<<{\sf kkkk}>> заданы как 


однородные, последовательность <<{\sf Kkkkk}>> будет расцениваться как один 


сегмент и~обрабатываться так же, как если бы в~обработку поступила 


последовательность из~7~идущих подряд символов~<<{\sf k}>> в~одном регистре.





      \smallskip


  


\noindent


\textbf{4.\ Приведение символьных последовательностей к~их 


формальным сжатым графематическим представлениям}


  


  Приведение осуществляется в~соответствии с~таблицей GA\_2, в~которой 


указаны интервалы. После сжатия текст будет иметь следующий вид:


  


    \noindent


    {\sf  %\small


      9L1,1Z2D1.2D1.4D1,1Z9L1-9L1Z9L1Z1D1Z9L1Z9L1Z9L1Z9L1Z1L1Z9L1Z1L1}


      


      \noindent


      {\sf Z9L1.}





      \smallskip


      


\noindent


\textbf{5.\ Поиск в~массиве формальных графематических 


представлений шаблонов формул по словарю и~определение типа 


получившегося токена}


  


  Для <<сборки>> сегментов в~один токен и~определения его типа применяется 


словарь GA\_3 составных сегментов. Если составная формула группы рядом 


стоящих сегментов есть в~словаре GA\_3, то такие сегменты подлежат 


объединению. Если нет, то группа сегментов уменьшается на~1 справа и~поиск 


продолжается.


  


  Вначале требуется определить максимальную длину формулы в~символах 


($F_{\max}$) и~вычислить на ее основе максимальное число сегментов, 


которые могут быть сгруппированы ($N_{\max}$):


     $$


     N_{\max}=\fr{F_{\max}}{2}\,.


     $$


  


  В данном примере $F_{\max}=10$ символов; следовательно, 


$N_{\max}\hm=5$~сегментов.


  


  Далее последовательно с~первого сегмента объединяются $N_{\max}$  


сегментов и~в~строку суммируются их формулы для получения искомой 


составной формулы, которую требуется проверить по словарю GA\_3 на 


возможность объединения.


  


  В данном примере для сегментов 1--5 составная формула =\;<<{\sf 9L1,1Z2D1}.>>.


  


  Если результат поиска составной формулы в~словаре GA\_3 отрицательный, 


число сегментов сокращается на~1~сегмент справа и~процесс поиска 


продолжается до тех пор, пока не будет найдена составная формула, или до тех 


пор, пока в~словаре GA\_3 не будет найдена формула оставшегося одного 


простейшего сегмента.


  


  Для анализируемой группы в~словаре GA\_3 находится только простейшая 


формула для первого сегмента~--- 9L. Тип такого токена~--- Word.


  


  Далее проверяются на совместимость следующие $N_{\max}$  сегментов, 


начиная с~сегмента, следующего за последним сегментом, вошедшим в~токен. 


Проверка продолжается до включения в~токены последнего сегмента текста. 


При этом делимитеры токенами не являются.


  


  В данном примере:


  \begin{itemize}


\item[\,]  сегменты 4--8 (составная строка классов <<{\sf DD.DD.DDDD}>>, составная 


формула <<{\sf 2D1.2D1.4D}>>) были найдены в~словаре GA\_3 и~объединены 


в~один токен типа Number, класс токена~12 (цифровой формат даты);


  \item[\,]


  сегменты 11--13 (составная строка классов <<{\sf kkk-kkkkk}>>, составная 


формула <<{\sf 9L1-9L}>>) были найдены в~словаре GA\_3 и~объединены в~один 


токен типа Word, класс токена~2 (сложносоставное слово);


  


  сегменты 3 и~10, содержащие пробелы, определились по словарю GA\_3 как 


делимитеры (формула <<{\sf 1Z}>>) и~в массив токенов не вошли.


\end{itemize}





  


\noindent


\textbf{6.\ Получение по имеющимся адресам сегментов исходного 


текста, разбитого на токены}





\section{Заключение}


  


  Широко распространенные способы токенизации не позволяют в~полной 


мере преобразовать символьное представление сложных технических текстов 


в~сис\-те\-му информационных объектов. Методы, основанные на правилах, не 


могут учесть все разнообразие сложных случаев, а нейросети реализуют только 


те ситуации, которые встречались в~обучающей выборке. Метод, описанный 


в~статье, учитывает указанные недостатки и~дает возможность успешно 


токенизировать сложные текстовые объекты благодаря применению 


функциональных шаблонов. К~их дополнительным преимуществам можно 


отнести возможность управления включением или выключением токенизации 


определенного вида токенов, добавления новых шаблонов по образцу текста 


и~получения на стадии формирования токенов их детальной классификации.








{\small\frenchspacing


{%\baselineskip=10.8pt


\addcontentsline{toc}{section}{Литература}


\begin{thebibliography}{9}  


  \bibitem{1-h}


  \Au{Ермакович М.\,В.} Автоматическое определение границ слова в~русском тексте  


с~по\-мощью комплекса лингвистических правил~/\!\!/ Компьютерная лингвистика 


и~интеллектуальные технологии: по мат-лам междунар. конф. <<Диалог 2017>>.~--- М.: РГГУ, 


2017. {\sf http://www.dialog-21.ru/media/3981/yermakovichmv.pdf}.


  \bibitem{2-h}


  \Au{Белоногов Г.\,Г., Калинин~Ю.\,П., Хорошилов~А.\,А.} Компьютерная лингвистика 


и~перспективные информационные технологии. Теория и~практика построения сис\-тем 


автоматической обработки текстовой информации.~--- М.: Русский мир, 2004. 248~с. 


  \bibitem{3-h}


  \Au{Гречачин В.\,А.} К вопросу о токенизации текста~/\!\!/ Международный  


на\-уч\-но-ис\-сле\-до\-ва\-тель\-ский ж., 2016. №\,6-4(48).С.~25--27.


 \end{thebibliography}


} }











\vspace*{-6pt}





\hfill{\small\textit{Поступила в~редакцию 15.09.22}}





%\vspace*{8pt}





%\hrule





%\vspace*{2pt}





\newpage





\vspace*{-28pt}





%\hrule





%\vspace*{8pt}





\def\tit{TOKENIZATION BASED ON~THE~METHOD OF~FUNCTIONAL~PATTERNS}











\def\titkol{Tokenization based on~the~method of~functional patterns}





\def\aut{Yu.\,V.~Nikitin$^1$, A.\,A.~Khoroshilov$^{1,2,3}$, and~A.\,E.~Makarova$^4$}





\def\autkol{Yu.\,V.~Nikitin, A.\,A.~Khoroshilov, and~A.\,E.~Makarova}





\titel{\tit}{\aut}{\autkol}{\titkol}





\vspace*{-8pt}





\noindent


   $^1$Federal Research Center ``Computer Science and Control'' of the Russian Academy\linebreak


$\hphantom{^1}$of Sciences, 44-2~Vavilov Str., Moscow 119333, Russian Federation


   


   \noindent


   $^2$Moscow State Aviation Institute (National Research University),  4~Volokolamskoe\linebreak


$\hphantom{^1}$Shosse, Moscow 125933, Russian Federation


   


   \noindent


   $^3$27th Central Research Institute of the Ministry of Defence of the Russian Federa-\linebreak 


$\hphantom{^1}$tion, 5,~1st~Khoroshevsky Passage, Moscow 123007, Russian Federation


   


   \noindent


   $^4$Scientific Industrial Joint Stock Company ``High Technology and Strategic Sys-\linebreak


$\hphantom{^1}$tems,'' 27-9~Elektrozavodskaya Str., Moscow 107023, Russian Federation





\def\leftfootline{\small{\textbf{\thepage}


\hfill Sistemy i Sredstva Informatiki~---


Systems and Means of Informatics 2022 vol~32 no\,4}


}%


 \def\rightfootline{\small{Sistemy i Sredstva Informatiki~---


 Systems and Means of Informatics 2022 vol~32 no\,4


\hfill \textbf{\thepage}}}





\vspace*{3pt}


  


  


  


  


  \Abste{The article proposes a new method of text tokenization based on the 


use of generalized functional templates. The method is based on the classification of 


Unicode characters in terms of their role in the formation of text elements and on the 


use of compound patterns from the generalized character classes. Widespread regular 


expressions are not used here. A~specific feature of the method is the use of a~sequence of 


characters as a~part of the interval template. The strengths of the method include 


successful tokenization of complex information objects (numbers, geographic 


coordinates, names of articles of engineering products, etc.), obtaining the detailed 


classification of tokens at the stage of their formation, the ability to turn on and off 


tokenization of a~certain type of tokens, as well as adding new templates according to 


the sample text for additional training of the system.}


  


  \KWE{tokenization; segmentation; graphematic analysis; computational 


linguistics; patterns; substitution; token}


  





\DOI{10.14357\!/\!08696527220406}  





%\vspace*{-10pt}








%\Ack


%\noindent








%\vspace*{-3pt}





\renewcommand{\bibname}{\protect\rmfamily References}


%\renewcommand{\bibname}{\large\protect\rm References}





{\small\frenchspacing


{ %\baselineskip=12pt


\addcontentsline{toc}{section}{References}


\begin{thebibliography}{9}


  \bibitem{1-h-1}


  \Aue{Ermakovich, M.\,V.} 2017. Avtomaticheskoe opredelenie granits slova 


v~russkom tekste s~po\-moshch'yu kompleksa lingvisticheskikh pravil [Automatic 


word boundary disambiguation on russian text using linguistic rules]. 


\textit{Computational Linguistics and Intellectual 


Technologies. Papers from the Annual Conference (International) ``Dialogue.''} 


Moscow: RGGU. Available at: {\sf  


https://www.dialog-21.ru/media/3981/\linebreak yermakovichmv.pdf} (accessed October~11, 


2022).


  \bibitem{2-h-1}


  \Aue{Belonogov, G.\,G., Yu.\,P.~Kalinin, and A.\,A.~Khoroshilov.} 2004. 


\textit{Komp'yuternaya lingvistika i~perspektivnye informatsionnye tekhnologii. 


Teoriya i~praktika po\-stro\-eniya sistem avtomaticheskoy obrabotki tekstovoy 


informatsii} [Computational linguistics and advanced information technologies. 


Theory and practice of building systems for automatic processing of text 


information]. Moscow: Russkiy mir. 248~p. 


  \bibitem{3-h-1}


  \Aue{Grechachin, V.\,A.} 2016. K~voprosu o tokenizatsii teksta [The issue of 


text tokenization]. \textit{Mezhdunarodnyy nauchno-issledovatel'skiy zh.} 


[International Research~J.] 6-4(48):25--27.


\end{thebibliography}


} }





\vspace*{-6pt}





\hfill{\small\textit{Received September 15, 2022}} 





%\vspace*{-8pt} 





  


  \Contr


  


  \noindent


  \textbf{Nikitin Yury V.} (b.\ 1977)~--- scientist, Federal Research Center 


``Computer Science and Control'' of the Russian Academy of Sciences, 44-2~Vavilov 


Str., Moscow 119333, Russian Federation; \mbox{yuri.v.nikitin@gmail.com}


  


  \vspace*{3pt}


  


  \noindent


  \textbf{Khoroshilov Aleksander A.} (b.\ 1952)~--- Doctor of Science in 


technology, professor, Moscow State Aviation Institute (National Research 


University), 4~Volokolamskoe Shosse, Moscow 125933, Russian Federation; leading 


scientist, Federal Research Center ``Computer Science and Control'' of the Russian 


Academy of Sciences, \mbox{44-2}~Vavilov Str., Moscow 119333, Russian Federation; 


senior scientist, 27th Central Research Institute of the Ministry of Defence of the 


Russian Federation, 5, 1st~Khoroshevsky Passage, Moscow 123007, Russian 


Federation; \mbox{khoroshilov@mail.ru}


  


  \vspace*{3pt}


  


  \noindent


  \textbf{Makarova Anna E.} (b.\ 1997)~--- mathematical software developer, 


Scientific Industrial Joint Stock Company ``High Technology and Strategic Systems,'' 


\mbox{27-9}~Elektrozavodskaya Str., Moscow 107023, Russian Federation; 


\mbox{anna20497@list.ru}


  


  


\label{end\stat}





\renewcommand{\bibname}{\protect\rm Литература} 





  


